\section{Background}

\subsection{LTN Example}
\label{appendix:example_sat_ltn}
Let us denote the predicate $\Friend$ as $\shortFriend$ for brevity.
In the expression $\shortFriend(\shortalice,\shortbob)$, let $\Gt(\shortalice)$ and $\Gt(\shortbob)$ be vector embeddings in $\R^m$.
A primitive approximation of the friendship relationship could be a cosine similarity function $\Gt(\shortFriend) : (\bx,\by)\mapsto \frac{\bx \cdot \by}{||\bx|| ||\by||} $.  
If $\Gt(\shortalice) = \begin{bsmallmatrix}3&2&0&5\end{bsmallmatrix}$ and $\Gt(\shortbob) = \begin{bsmallmatrix}1&0&0&4\end{bsmallmatrix}$, we have 
$\Gt(\shortFriend(\shortalice,\shortbob)) = 0.905$; that is, a high truth degree. 

Of course, stating that people are friends if they are similar is primitive.
In a real-case scenario, the friendship relationship would likely be approximated by a parametric function, such as a neural network, 
and trained based on constraints in a loss function.

Consider the formula $\phi = \lnot \shortFriend(\shortalice, \shortbob) \lor \shortFriend(\shortbob, \shortalice)$
and a knowledgebase that contains this unique formula $\K = \{ \phi \}$.
The formula states that if $\alice$ is a friend of $\bob$, then $\bob$ is a friend of $\alice$.
    \footnote{$\phi$ is semantically equivalent to $\shortFriend(\shortalice, \shortbob) \imp \shortFriend(\shortbob, \shortalice)$ 
            if we use a material implication defined as $p \imp q \equiv \lnot p \lor q$.}
Let the grounding for $\Gt(\shortFriend(\shortalice, \shortbob))$ depend 
on a trainable neural network for the friendship relation
and a set of features for $\alice$ and $\bob$.
To update $\btheta$ via gradient descent steps, we calculate 
$\frac{\partial \Gt(\phi)}{\partial \Gt(\shortFriend(\shortalice,\shortbob))}$ 
and 
$\frac{\partial \Gt(\phi)}{\partial \Gt(\shortFriend(\shortbob,\shortalice))}$.
Using the operators $\Ns(x) = 1-x$ and $\Sprod(x,y)=x+y-xy$, we get:
%
\begin{align}
    \Gt(\phi) & =  \Sprod(\Ns(\Gt(\shortFriend(\shortalice,\shortbob))), \Gt(\shortFriend(\shortbob,\shortalice))) \\
            & = 1 - \Gt(\shortFriend(\shortalice,\shortbob)) + \Gt(\shortFriend(\shortalice,\shortbob)) \Gt(\shortFriend(\shortbob,\shortalice))
\end{align}
And the partial derivatives:
\begin{align}
    \label{eq:partial_friendba} \frac{\partial \Gt(\phi)}{\partial \Gt(\shortFriend(\shortbob,\shortalice))} & = \Gt(\shortFriend(\shortalice,\shortbob)) \\
    \label{eq:partial_friendab} \frac{\partial \Gt(\phi)}{\partial \Gt(\shortFriend(\shortalice,\shortbob))} & = -1+\Gt(\shortFriend(\shortbob,\shortalice))
\end{align}%

Equations \eqref{eq:partial_friendab} and \eqref{eq:partial_friendba} give us interesting insights on the power of LTN.
When maximizing the satisfiability of the formula, if $\shortFriend(\shortalice,\shortbob)$ is high,
then $\frac{\partial \Gt(\phi)}{\partial \Gt(\shortFriend(\shortbob,\shortalice))}$ is high.
Intuitively, if $\alice$ being friend with $\bob$ has a high truth value, LTN will tend to increase the truth value of $\bob$ being friend with $\alice$.
Alternatively, if $\shortFriend(\shortbob,\shortalice)$ is low, $\frac{\partial \Gt(\phi)}{\partial \Gt(\shortFriend(\shortalice,\shortbob))}$ is close to $-1$.
That means that if $\bob$ is not considered friend with $\alice$, LTN will tend to decrease the truth value of $\alice$ being friend with $\bob$.

These are different scenarios and ways in which the framework pushes parametric groundings to verify logical constraints.
One can easily imagine how the LTN loss can be used as an additional loss term when training neural networks or embeddings to find a balanced optimum that also satisfies a knowledgebase.


\subsection{Product Real Logic and Prenex Normal Form}
\label{appendix:prodrl_and_pnf}
The study conducted by \cite{van_krieken_analyzing_2022} evaluated a range of operators for differentiability and found Product Real Logic to be the current state-of-the-art operator semantics for differentiable fuzzy logics.
This set of semantics uses the product t-norm $\Tprod(x,y)=xy$, its dual t-conorm $\Sprod(x,y)=x+y-xy$, the standard negation $\Ns(x)=1-x$, and the material implication.
The universal quantifier uses the log-product aggregator $(\log \circ \Atprod)(x_1,\dotsc,x_n) = \sum_{i=1}^{n} \log(x_i)$,
and the existential quantifier uses a smooth maximum.

A limitation of this set of semantics is that it combines operators in both the standard space and the logarithm space, making it challenging to handle certain formulas. 
For example, a formula such as $(\forall u P(u)) \lor (\forall v Q(v))$ cannot be grounded as is, as the $\forall$ operator outputs a log truth degree while the $\lor$ operator expects a normal truth degree. 
One potential solution is to transform the formula into Prenex Normal Form (PNF) $\forall u \forall v (P(u) \lor Q(v))$, but this adds considerable complexity as we now need to ground combinations of individuals from $u$ and $v$. 
PNF formulas that contain universal quantifiers within the scope of existential quantifiers, such as $\forall u \exists v \forall w P(u,v,w)$, are even more difficult to handle. 
\logLTN, on the other hand, is simpler to work with in these cases.

\subsection{Simplification of Log Sigmoid and Log Softmax}
\label{appendix:logsigmoid}
Given the sigmoid function $\sigmoid(x) = \frac{1}{1+e^{-x}}$, for large negative values of $x$, we have $\log(\sigmoid(x)) = \log(\frac{1}{1+N}) = \log(1)-\log(1+N) \approx -N$ where $N$ is a large number. 
The derivative is also simple. Given that $\partialderivative{\sigmoid(x)}{x} = \sigmoid(x) (1-\sigmoid(x))$, we have:
\begin{equation}
    \partialderivative{\log(\sigmoid(x))}{x} = 1-\sigmoid(x)
\end{equation}
Similarly, given that the softmax function $\softmax(\mathbf{z})_i = \frac{e^{z_i}}{\sum_{j=1}^{K} e^{z_j}}$ over a vector $\mathbf{z}$ of $K$ values, $i=1,\dots,K$, 
has the derivatives $\partialderivative{\softmax(\mathbf{z})_i}{\mathbf{z}_j} = \softmax(\mathbf{z})_i (\delta_{ij}- \softmax(\mathbf{z})_j)$, we have:
\begin{equation}
    \partialderivative{\log(\softmax(\mathbf{z})_i)}{\mathbf{z}_j} = \delta_{ij}- \softmax(\mathbf{z})_j
\end{equation}
where $\delta_{ij} = \begin{cases}
        1 \ \ i=j\\
        0 \ \ i \neq j 
    \end{cases}$. 



\section{Experiments}
\subsection{Semantic PASCAL-Part Dataset}
\label{appendix:pascalpart}
The semantic PASCAL-Part dataset is a simplified version of the PASCAL-Part dataset introduced by \cite{pascalpart_chen_2014}. 
The goal is to train a type classifier $\istype(x,\bottle)$, $\istype(x,\captype)$, etc., that predicts the type of an object within a bounding box $x$,
 and to train a relation predictor $\ispartof(x,y)$ that determines if one bounding box $x$ is part of another bounding box $y$.

\subsubsection{Constraints}
Training is guided by three constraints based on ground truth examples,
one for labeled type examples, one for pairs of positive examples for $\ispartof$, and one for pairs of negative examples. 
Note that the negative pairs are always sampled in bounding boxes belonging to the same image.
\begin{gather}
    \forall \mathrm{diag}(x_\mathsf{label}, \mathsf{label})\ \istype(x_\mathsf{label}, \mathsf{label}) \\
    \forall \mathsf{pairs_+}\ \ispartof(\mathsf{pairs_+[0]}, \mathsf{pairs_+[1]}) \\
    \forall \mathsf{pairs_-}\ \lnot \ispartof(\mathsf{pairs_-[0]}, \mathsf{pairs_-[1]}) \\
\end{gather}
Where $\mathrm{diag}(x_\mathsf{label}, \mathsf{label})$ is a special quantification that aggregates only arranged pairs of bounding boxes and their labels,
$\mathsf{pairs_+}$ is a batch of positive examples of $\ispartof$, $\mathsf{pairs_-}$ is a batch of negative examples of $\ispartof$.

There are two constraints stating that $\ispartof$ is antisymmetric and antireflexive.
\begin{gather}
    \forall \mathsf{pairs}\ \lnot \big( \ispartof(\mathsf{pairs[0]}, \mathsf{pairs[1]}) \land \ispartof(\mathsf{pairs[1]}, \mathsf{pairs[0]}) \big) \\
    \forall x\ \lnot \ispartof(x, x)
\end{gather}

Finally, and most importantly, there are mereological constraints that the types and their meanings.
The mereological constraints are based on the ontologies in Table \ref{tab:ontologies} rearranged in the shape of \eqref{eq:constraint_mereo1_appendix} and \eqref{eq:constraint_mereo2_appendix}.
\begin{gather}
    \forall x,y \ \istype(x,\bottle) \land \ispartof(y,x) \imp (\istype(y,\captype) \lor \istype(x,\body)) \label{eq:constraint_mereo1_appendix} \\
    \forall x,y \ \istype(x,\captype) \land \ispartof(x,y) \imp \istype(y,\bottle) \label{eq:constraint_mereo2_appendix} 
\end{gather}

\subsubsection{Features}
This setup has been previously implemented by \cite{donadello_ltnimage_2017,van_krieken_semi-supervised_2019}. 
In these previous works, the bounding boxes were grounded with the object class predictions produced by an object detector trained on the PASCAL-Part dataset. 
This means that LTN was only used to correct the predictions of the detector. 
In contrast, we increased the difficulty by grounding the bounding boxes with a latent vector of 1024 features, which is produced by an intermediate layer of the FasterRCNN~\cite{Ren_fasterrcnn_2017}.
This means that LTN must also learn the final layers of the object classifier.
In addition, we included the coordinates of each bounding box and their overlapping ratio when grounding the pairs.

For a work that trains an object detector architecture end-to-end on all types using LTN, refer to \cite{manigrasso_faster-ltn_2021}.

\subsubsection{Violation metrics}
In addition to the standard accuracy metrics for each predictor (PR AUC for the part-of predictor and balanced accuracy for the type predictor), we also assess their combined performance by measuring \emph{violations} of the mereological constraints.
It's important to note that not all misclassifications are equal. For instance, if a model predicts a cap inside a plant, it suggests that the system has learned less from prior knowledge compared to a model that misclassifies the cap in a different context, such as a wrong bottle.

This concept is visually represented in Figure \ref{fig:examples_violations}. Our results demonstrate that not only does \logLTN\ exhibit significantly fewer violations, but it also reaches this outcome much faster compared to other baseline approaches, as shown in Figure \ref{fig:pascalpart_violations}.

\begin{figure}
    \centering
    \begin{subfigure}{0.45\textwidth}
      \centering
      \includegraphics[width=0.8\linewidth]{figures/nonviolation2.png}
      \caption{Nonviolation (a cap port of a bottle)}
      \hfill
    \end{subfigure}
    ~
    \begin{subfigure}{0.45\textwidth}
      \centering
      \includegraphics[width=0.8\linewidth]{figures/violation1.png}
      \caption{Violation (a plant part of another plant)}
    \end{subfigure}
    
    \caption{Examples of false positives on Semantic PASCAL-Part. Green is predicted in blue.}
    \label{fig:examples_violations}
  \end{figure}
  


\begin{figure}
    \centering
    \includegraphics[width=0.5\linewidth]{figures/pascalpart_violations_results.pdf}
    \caption{Violations over training time. The first 100 training steps are used for pretraining without the mereological constraints.}
    \label{fig:pascalpart_violations}
\end{figure}


\begin{table}
    \begin{tabular}{l|l}
        \toprule
        Whole & Parts \\
        \midrule
        aeroplane & artifact\_wing, body, engine, stern, wheel \\
        bicycle & chain\_wheel, handlebar, headlight, saddle, wheel \\
        bird & animal\_wing, beak, tail, eye, head, leg, neck, torso \\
        bottle & body, cap \\
        bus & bodywork, door, headlight, license\_plate, mirror, wheel, window \\
        car & bodywork, door, headlight, license\_plate, mirror, wheel, window \\
        cat & ear, tail, eye, head, leg, neck, torso \\
        cow & ear, horn, muzzle, tail, eye, head, leg, neck, torso \\
        dog & ear, muzzle, nose, tail, eye, head, leg, neck, torso \\
        horse & ear, hoof, muzzle, tail, eye, head, leg, neck, torso \\
        motorbike & handlebar, headlight, saddle, wheel \\
        person & arm, ear, ebrow, foot, hair, hand, mouth, nose, eye, head, leg, neck, torso \\
        pottedplant & plant, pot \\
        sheep & ear, horn, muzzle, tail, eye, head, leg, neck, torso \\
        train & coach, headlight, locomotive \\
        tvmonitor & screen \\
        boat & \\
        chair & \\
        sofa & \\
        diningtable & \\
        \bottomrule
    \end{tabular}
\caption{Full ontologies in Semantic PASCAL-Part}
\label{tab:ontologies}        
\end{table}


\subsection{Implementation Details}
\label{appendix:implementation}

logLTN is made available as a subpackage of the LTN library.\footnote{\url{https://github.com/logictensornetworks/logictensornetworks}}
Table \ref{tab:appendix_models} details the neural models used in each experimental task and Table \ref{tab:appendix_operator_params} details the hyperparameters of the baselines.
The Adam optimizer is trained with a learning rate of $0.002$ in the clustering task and a learning rate of $0.001$ for the MNISTAdd and Semantic Image Interpretation (SII) task.
We run our experiments on a machine equipped with a Tesla T4 GPU.

\begin{table}
    \begin{center}
    \begin{tabular}{llll}
        \toprule
        Baseline  & Operator & Parameter & Schedule \\
        \midrule
        logLTN & $\exists: \LME$ & $\alpha$ & Linear: $[1,4]$\\
        Prod RL & $\exists: \mathrm{pM}$ & $p$ &  Linear: $[1,6]$\\
        Stable RL & $\exists: \mathrm{pM}$ & $p$ &  Linear: $[1,6]$\\
        \bottomrule
    \end{tabular}
    \end{center}
    
    \begin{tablenotes}
        \small
        \item Linear $[a,b]$: the parameter increases linearly from $a$ to $b$ over the training steps,
        \item $\mathrm{pM}(x_1, \dots, x_n \mid p) = (\frac{1}{n} \sum_{i=1}^{n} x_i^p)^{\frac{1}{p}}$,
        \item $\mathrm{pME}(x_1, \dots, x_n \mid p) = 1 - (\frac{1}{n} \sum_{i=1}^{n} (1- x_i)^p)^{\frac{1}{p}}$.
      \end{tablenotes}
    \caption{Overview of the hyperparameters used in the experiments for each fuzzy operator configuration.}
    \label{tab:appendix_operator_params}
\end{table}

\newcommand{\Conv}{\mathrm{Conv}}
\newcommand{\Dense}{\mathrm{Dense}}

\begin{table*}[b]
    \centering
    \begin{tabular}{llll}
        \toprule
        Task  & Predicate & Model & Output layer \\
        \midrule
        Clustering & $C(x,c)$ & $\mathrm{Dense}(16)^\ast$, $\mathrm{Dense}(16)^\ast$, $\mathrm{Dense}(16)$ & $\mathrm{Softmax}$ \\
        MNISTAdd & $\isdigit(x,d)$ & \makecell[tl]{$\Conv(6,5)^\ast$, $\mathrm{MP}(2,2)$, $\Conv(16,5)^\ast$, $\mathrm{MP}(2,2)$,\\
                 $\Dense(100)^\ast$, $\Dense(84)^\ast$, $\Dense(10)$} & $\mathrm{Softmax}$\\
        SII & $\mathsf{type}(x,t)$ & \makecell[tl]{$\mathrm{Dense}(512)^\ast$, $\mathrm{Dense}(256)^\ast$, $\mathrm{Dense}(256)^\ast$,\\ 
                $\mathrm{Dense}(128)^\ast$, $\mathrm{Dense}(128)$} & $\mathrm{Softmax}$ \\
        SII & $\mathsf{partof}(x,y)$ & \makecell[tl]{$\mathrm{Concat}(\mathsf{type}_{\mathrm{Model}}(x),\mathsf{type}_{\mathrm{Model}}(y))$, $\mathrm{Dense}(512)^\ast$, \\ 
                $\mathrm{Dense}(256)^\ast$, $\mathrm{Dense}(256)^\ast$, $\mathrm{Dense}(128)^\ast$, $\mathrm{Dense}(128)$}      &  $\mathrm{Sigmoid}$ \\
        \bottomrule
    \end{tabular}
    \begin{tablenotes}
        \small
        \item $\ast$: layer ends with an $\mathsf{elu}$ activation,
        \item $\mathrm{Dense}(k)$: linear layer with $k$ units,
        \item $\Conv(f,k)$ : 2D convolution layer with $f$ filters and a kernel of size $k$,
        \item $\mathrm{MP}(w,h)$ : max pooling operation with a $w \times h$ pooling window.
      \end{tablenotes}
    \caption{Overview of the neural architectures used in each task.}
    \label{tab:appendix_models}
\end{table*}

\section{Theory}
\subsection{Log-Negations}
\subsubsection{Proof of Theorem \ref{thm:log_not_sigmoid}}
\label{appendix:proof_log_not_sigmoid}
\lognotsigmoid*
\begin{proof}
    \begin{align*}
        (\log \circ \Ns) (x) &= \log(1-S(x)) = \log(S(x))+\log(\frac{1-S(x)}{S(x)})\\
        & =\log(S(x)) + \log(\frac{1 - \frac{e^x}{e^x+1}}{ \frac{e^x}{e^x+1}}) \\
        & = \log(S(x)) + \log(\frac{e^x + 1 - e^x}{e^x})\\
        & = \log(S(x)) - x \qedhere
    \end{align*}
\end{proof}
    
\subsubsection{Proof of Theorem \ref{thm:log_not_softmax}}
\label{appendix:proof_log_not_softmax}
\lognotsoftmax*
\begin{proof}
    \begin{align*}
        (\log \circ \Ns) (\softmax(\mathbf{z})_i)
         &= \log(1-\softmax(\mathbf{z})_i) \\
         & = \log(\softmax(\mathbf{z})_i) + \log(\frac{1-\softmax(\mathbf{z})_i}{\softmax(\mathbf{z})_i})\\
         &=  \log(\softmax(\mathbf{z})_i) + \log(\frac{1 - \frac{e^{z_i}}{\sum_{j=1}^{K} e^{z_j}}}{\frac{e^{z_i}}{\sum_{j=1}^{K} e^{z_j}}})\\
        & =  \log(\softmax(\mathbf{z})_i) + \log(\frac{\sum_{j=1}^{K} e^{z_j} - e^{z_i}}{e^{z_i}}) \\
        &= \log(\softmax(\mathbf{z})_i) + \log(\sum_{\substack{j=1\\j\neq i}}^K e^{z_j}) - z_i \qedhere
    \end{align*}
\end{proof}

\subsection{Relaxation of the Disjunction}
\subsubsection{Bounds of LogSumExp}
\label{appendix:bounds_LSE_lower}
We start from Equation \eqref{eq:bounds_LSE_traditional} and substract $-\frac{\log(n)}{a}$ in all parts of the inequality:
\begin{equation}
    \label{eq:appendix_LSE_lb_1}
    \max(\mathbf{x}) - \frac{\log(n)}{\alpha} \leq \LSE(\mathbf{x} \mid \alpha, C) - \frac{\log(n)}{\alpha} \leq \max(\mathbf{x})
\end{equation}

And:
\begin{align}
    \LSE(\mathbf{x} \mid \alpha, C) - \frac{\log(n)}{\alpha} &= \frac{1}{\alpha} \Bigg( C + \log(\sum_{i=1}^n e^{\alpha x_i - C}) - \log(n)\Bigg) \\
            & \quad = \frac{1}{\alpha} \left( C + \log(\frac{\sum_{i=1}^n e^{\alpha x_i-C}}{n} ) \right) \\
   \label{eq:appendix_LSE_lb_2}         & \quad = \LME(\mathbf{x} \mid \alpha, C)
\end{align} 

Given \eqref{eq:appendix_LSE_lb_1} and \eqref{eq:appendix_LSE_lb_2}, we obtain the inequality for $\LME$:
\begin{equation}
    \max(\mathbf{x}) - \frac{\log(n)}{\alpha} \leq \LME(\mathbf{x} \mid \alpha, C) \leq \max(\mathbf{x})
\end{equation}



\subsection{De Morgan's Inequalities}
\subsubsection{Proof of Theorem \ref{thm:log_demorgans}}
\label{appendix:proof_nnfdemorgans}
\logdemorgans*

\begin{proof}
    Let $\G(P) = x$, $\G(Q) = y$, and $\G(P(u)) = [x_1, \dots, x_n]$, where $x, y, x_1, \dots, x_n \in [0,1]$.
    Grounding the operators with the definitions of Section \ref{s:logltn_semantics}, we need to prove:
    \begin{align}
        \Ns(\Tprod(x,y)) &\geq \Smax(\Ns(x),\Ns(y)) \label{eq:demorgan_tprod}\\
        \Ns(\Smax(x,y)) &\geq \Tprod(\Ns(x),\Ns(y)) \label{eq:demorgan_smax}\\
        \Ns(\Atprod(x_1,\cdots,x_n)) &\geq \Asmax(\Ns(x_1),\dots,\Ns(x_n)) \label{eq:demorgan_atprod}\\
        \Ns(\Asmax(x_1,\cdots,x_n)) &\geq \Atprod(\Ns(x_1),\dots,\Ns(x_n)) \label{eq:demorgan_asmax}
    \end{align}
    \eqref{eq:demorgan_tprod} and \eqref{eq:demorgan_smax} are specializations of the other two equations by working with $n=2$ values.
    We focus on proving \eqref{eq:demorgan_atprod}.
    Indeed, we can easily retrieve \eqref{eq:demorgan_asmax} from \eqref{eq:demorgan_atprod}.
    \paragraph{Equivalence of \eqref{eq:demorgan_atprod} and \eqref{eq:demorgan_asmax}}
    Posing $x_i'=1-x_i$ for $i=1\dots n$, we have:
    \begin{align}
        & \quad \Ns(\Atprod(x_1,\cdots,x_n)) \geq \Asmax(\Ns(x_1),\dots,\Ns(x_n))\\
        \iff & \quad \Ns(\Ns(\Atprod(x_1,\cdots,x_n))) \leq \Ns(\Asmax(\Ns(x_1),\dots,\Ns(x_n))) \\
        \iff & \quad \Atprod(x_1,\cdots,x_n) \leq \Ns(\Asmax(\Ns(x_1),\dots,\Ns(x_n))) \\
        \iff & \quad \Atprod(\Ns(x_1'),\cdots,\Ns(x_n')) \leq \Ns(\Asmax(x_1',\dots,x_n'))\\
        \iff & \quad \Ns(\Asmax(x_1',\dots,x_n')) \geq \Atprod(\Ns(x_1'),\cdots,\Ns(x_n'))
    \end{align}

    \paragraph{Proof of \eqref{eq:demorgan_atprod}}
    In the left-hand side of the inequality, we have:
    \begin{equation}
        \Ns(\Atprod(x_1,\cdots,x_n)) = 1-\prod_{i=1}^{n} x_i
    \end{equation}
    and in the right-hand side:
    \begin{equation}
        \Asmax(\Ns(x_1),\dots,\Ns(x_n)) = \max_{i=1}^n (1-x_i) = 1-\min_{i=1}^n (x_i)
    \end{equation}
    Replacing them in the original inequality, and denoting $j = \argmin_{i=1}^n (x_i)$, we obtain:
    \begin{align}
        & \quad 1-\prod_{i=1}^{n} x_i \geq 1-\min_{i=1}^n (x_i) \\
        \iff & \quad \prod_{i=1}^{n} x_i \leq \min_{i=1}^n (x_i) \\
        \iff & \quad x_j \prod_{\substack{i=1\\i\neq j}}^n x_i \leq x_j \\
        \iff & \quad \prod_{\substack{i=1\\i\neq j}}^n x_i \leq 1
    \end{align}
    Which is true because all $x_i \in [0,1]$.
\end{proof}

\subsubsection{Tightness of the bounds}
\label{appendix:bound_tightness}
We provide an analysis of the tightness of the bounds of the De Morgan's inequalities.
We measure the tightness for the quantifier variants of the inequalities, as this generalizes to the case $n=2$.
\paragraph{Tightness of \eqref{eq:demorgan_atprod}}
Let us first characterize the maximum value of the bound.
We are interesting in finding the values $(x_1^*,\dots,x_n^*)$ that maximize the difference of the two members in \eqref{eq:demorgan_atprod}.
\begin{align}
    (x_1^*,\dots,x_n^*) = \argmax_{(x_1,\dots,x_n)}  \Delta_\land(x_1,\dots,x_n)
\end{align}
with
\begin{align}
     \Delta_\land(x_1,\dots,x_n) &= \Ns(\Atprod(x_1,\cdots,x_n)) - \Asmax(\Ns(x_1),\dots,\Ns(x_n)) \\
            &= (1-\prod_{i=1}^{n} x_i) - (1-\min_{i=1}^n(x_i) ) = \min_{i=1}^n(x_i) - \prod_{i=1}^n x_i 
\end{align}
Let $j = \argmin_{i=1}^n (x_i)$.
\begin{align}
    \Delta_\land(x_1,\dots,x_n) = x_j - \prod_{i=1}^n x_i =  x_j (1-\prod_{\substack{i=1\\i\neq j}}^n x_i)
\end{align}
For any set of values, we have $\Delta_\land(x_j,\dots,x_j) \geq \Delta_\land(x_1,\dots,x_n)$, 
as $(1-x_j^{n-1}) \geq (1-\prod_{\substack{i=1\\i\neq j}}^n x_i)$ given that $x_j = \min_{i=1}^n (x_i)$.
Therefore, $(x_1^*,\dots,x_n^*)$ is actually a tuple of the same value taken $n$ times,
and we reduce the search to:
\begin{align}
    x^* = \argmax_{x}  \Delta_\land(x) = \argmax_x x - x^n
\end{align}
Given that $x \in [0,1]$, we find by first and second order derivative analysis that $\Delta_\land(x)$ is concave on the whole domain and has a single maximum at:
\begin{align}
    \label{eq:demorgan_and_max_bound}
    x^* & = n^{-(\frac{1}{n-1})} 
\end{align}

For $n=2$, that is, when applying the De Morgan's law to a simple conjunction of two terms $x_1$ and $x_2$, the bound of the inequality is maximal when $x_1=x_2=0.5$, giving $\Delta_\land = 0.25$.
However, on average, the bound is smaller. 
By sampling 10e4 $\times$ 10e4 points linearly on the domain $[0,1] \times [0,1]$, we find an average $\Delta_\land = 0.083167$. 
We visualize the bound of the inequality in Figure \ref{fig:demorgan_bounds}.
It is zero when any $x_i=0$ or $x_i=1$, in which case the De Morgan's laws are verified.

For larger $n$, that is, when applying the quantifier equivalent of the De Morgan's law, the maximal value of the bound becomes larger.
For example, with $n=8$, the bound is maximal for $x_1=\dots=x_8\approx 0.743$, with $\Delta_\land \approx 0.650$.
However, on average, the bound stays small.
By sampling 10e8 points linearly in the domain $[0,1]^8$, we have an average bound of $0.07135$.

\begin{figure}
    \centering
    \includegraphics[width=0.8\textwidth]{figures/demorgan_bound_2d_and.pdf}
    \caption{De Morgan's inequality bound $-\Delta_\land(x_1,x_2)$}
    \label{fig:demorgan_bounds}
\end{figure}



\paragraph{Tightness of \eqref{eq:demorgan_asmax}}
As we have shown in Section~\ref{appendix:proof_nnfdemorgans}, \eqref{eq:demorgan_asmax} is equivalent to \eqref{eq:demorgan_atprod} by replacing the input values with $x_i'=1-x_i$ for $i=1\dots n$.
Therefore, the maximum value of the bound is the same, except that it happens on a complement set of value.
We have:
\begin{align}
    x^* & = \argmax_{x}  \Delta_\lor(x) \\ 
        &= \argmax_{x} \Ns(\Asmax(x,\cdots,x)) - \Atprod(\Ns(x),\dots,\Ns(x)) \\
            &= \argmax_x 1 - x - (1-x)^{n} \\
            \label{eq:demorgan_or_max_bound}
            &= 1 - n^{-(\frac{1}{n-1})} 
\end{align}

For $n=2$, when applying the De Morgan's law to a simple disjunction of two terms $x_1$ and $x_2$, $\Delta_\lor$ is maximal when $x_1=x_2=0.5$.
We find again that the peak is $\Delta_\lor = 0.25$, and by sampling 10e4 $\times$ 10e4 points linearly on the domain $[0,1] \times [0,1]$, we find an average $\Delta_\lor =  0.083167$. 

For $n=8$, this time the bound is maximal when $x_1=\dots=x_8\approx 0.257$, with $\Delta_\lor \approx 0.650$.
Sampling 10e8 points linearly in $[0,1]^8$, we still get an average bound of $0.07135$.


\subsection{Common Fuzzy Properties for \logLTN}
We provide an overview of common fuzzy properties that are verified by the operator configuration in Table \ref{tab:fuzzy_properties}. 
The distributivity of $\land$ over $\lor$ for \logLTN\ is the only new property and can be demonstrated easily.
Let $\G(P)=x$, $\G(Q)=y$, and $\G(R)=z$ be the grounding of three predicates. We have:
\begin{align}
    & \G(P \land (Q \lor R)) = \G((P \land Q) \lor (P \land R)) \\
    \iff & \qquad  x \max(y,z) = \max(xy, xz)
\end{align}


\begin{table}
    \centering
    \begin{tabular}{l|c|c|c}
        \toprule
        Property                                            & \logLTN       & Prod RL       & Stable RL \\
        \midrule
        Commutativity of $\land$, $\lor$                    & \checkmark    & \checkmark    &  \checkmark \\
        Associativity of $\land$, $\lor$                    & \checkmark    & \checkmark    &  \checkmark \\
        De Morgan's laws for $\land$ and $\lor$             &               & \checkmark    &  \checkmark \\
        Material Implication                                & \checkmark    & \checkmark    & \checkmark  \\
        Distributivity of $\land$ over $\lor$               & \checkmark    &               &             \\
        Distributivity of $\lor$ over $\land$               &               &               &             \\
        Double negation, i.e.\ $\G(\lnot \lnot p) = \G(p)$                 & \checkmark    & \checkmark    & \checkmark  \\
        Law of non-excluded middle, i.e.\ $\G(p \land \lnot p) = 0$                          &               &               &             \\
        Law of non-contradiction, i.e.\ $\G(p \lor \lnot p) = 1$                           &               &               &             \\
        Conjunction elimination, i.e.\ $\G(p \land p) \leq \G(p)$ & \checkmark & \checkmark & \checkmark \\
        Disjunction amplification, i.e.\ $\G(p \lor p) \geq \G(p)$ & \checkmark & \checkmark & \checkmark \\
        $\forall$ defined as a generalization of $\land$    & \checkmark    & \checkmark    &             \\
        $\exists$ defined as a generalization of $\lor$     & \checkmark    &               &             \\
        De Morgan's laws for $\forall$ and $\exists$        &               &               &             \\
        
        \bottomrule
    \end{tabular}
    \caption{Fuzzy properties for \logLTN, Product Real Logic (Prod RL), and Stable Product Real Logic (Stable RL).}
    \label{tab:fuzzy_properties}
\end{table}

